% The two phases of infection --- plan \S3.3, "\S8".  Opens Part B.
% Sources: 10/02 entire.
% C7: 10/02:57 prints p_0 with no derivation and no declaration that its
%     letters are CH4/CH5's own.  The formula is correct; the work here is the
%     derivation, the declaration and the closed form.
% A4: the geometric sum closed, \cref{var:eq:p0closed}.
% Renames (plan \S7.2): phase-2 rates \beta,\nu -> \lambda_2,\mu_2; phase-1
% rates \lambda,\mu,\delta -> \lambda_1,\mu_1,\delta_1; and the two constants
% of the decaying death rate, both printed as "a" at 10/02:14, written
% \mu_\infty and \mu_0 -- see the NEEDS-AUTHOR marker below, since a is CH4's
% larger root and is used in that sense in this very section.
% The source at 10/02:63 prints E(Z) = 2(1 - b + sum); the sign of the sum is
% wrong there, since p_0 = b + sum.  \cref{var:eq:EZ} carries the correct one.

\section{The two phases of infection}
\label{var:sec:twophase}

In every model so far a rate has moved because the process changed its own
circumstances: it used up a capacity, hardened a competitor, depleted a
potential, infected its neighbours, or manufactured its own advantage. This
last model is not like that. Here the rate changes because the host does. An
infection proceeds for a while under one set of conditions, the host
recognises it, and from that point the conditions are different ones. That
difference pays for itself at once: each phase is a constant-rate process of
a kind already solved, and the whole problem reduces to describing the state
handed across the join. This is the section in which the mathematics closes
again, and \cref{var:sec:pattern} takes up what exactly buys that.

\subsection{The biology, and the shape it dictates}
\label{var:sec:twophasebiology}

The pathogenesis of \yp{} has two distinct phases, the pre-inflammatory and
the post-inflammatory. In the first the bacterium replicates mostly
intracellularly, within macrophages
\cite{PujolBliska2003,ConnorEtAl2018}. As the body recognises the alien
presence and mounts a neutrophil response, it switches to extracellular
replication \cite{BosioEtAl2005,PetersEtAl2013}. That the bacterium does not
simply begin extracellularly is itself the evidence that living within
macrophages carries an advantage, and the simplest model of that advantage is
a higher replication rate during the first phase.

The switch from intracellular to extracellular replication is coincident
with the neutrophil response in the histology
\cite{BosioEtAl2005,PetersEtAl2013}; whether the neutrophils drive the
switch, rather than arriving at the same time, is not settled by those
studies and is not needed below.

\begin{remark}[A decaying extracellular death rate]
\label{var:rem:mudecay}
Inside the cell there are fewer of the extracellular mechanisms of innate
immunity --- complement, and bacteriophages. One reading, the least supported
claim in this section, is that the bacterium adapts to host-specific phages
within the handful of replication cycles a single infection lasts, so that
the death rate in the second phase decays rather than holding constant. That
would give
\begin{equation}
\mu_2(t)=\mu_\infty+\mu_0\,\ee^{-t/T},
\label{var:eq:mudecay}
\end{equation}
with $T$ the decay timescale, taken as the unit of time in any numerical
illustration so that the exponential is $e^{-t}$. The two constants are
independent; nothing requires the initial death rate to be twice the
asymptotic one. With this rate the second phase is subcritical for a while
and then not, and the size of the population handed over from the first phase
is decisive. The analysed model below does not use it.
\end{remark}

The same two-phase structure is argued from the other direction in
\cref{path:sec:ypestis}, where the phases are distinguished by the density of
the killer-cell field and the switch is placed where a removal budget
$\vartheta^{\ast}$ is crossed. That account asks why the switch happens;
this one asks what composing the two phases does to the probability that the
infection establishes. Neither is derived from the other and they are not
duplicates, but they concern one organism and one transition, and
\cref{var:sec:connections} returns to what reconciling them would take.

\subsection{Phase composition}
\label{var:sec:composition}

Take the first phase to be a birth--death--catastrophe process with rates
$\lambda_1,\mu_1,\delta_1$ and $\lambda_1\gg\mu_1$, so that it strongly
favours replication, running until the catastrophe --- the rupture of the
host cell --- releases its load. Take the second to be a linear birth--death
process with rates $\lambda_2,\mu_2$, only slightly supercritical, so that
$\lambda_2>\mu_2$ but $\lambda_2\approx\mu_2$. Write $t_c$ for the switch.

The biology of \cref{var:sec:twophasebiology} puts the switch at the
neutrophil response, and the model puts it at the rupture of the host cell.
These are treated here as one event: the rupture is what ends the
intracellular phase, and the immune response is what the surrounding argument
supposes provokes it. So $t_c$ is not a fixed time chosen from outside but the
catastrophe time of the first phase, and \cref{var:eq:phasepgf} is read
conditionally on it. That is the honest reading, and it is also the one that
explains why this section closes where \cref{var:sec:rrrr}, whose catastrophe
rate is state-dependent in the same way, does not: the burst-size law
\cref{dist:eq:burstlaw} has already integrated the rupture time out, and hands
across a distribution over $\Kb$ that no longer mentions when the rupture
came.

\emph{Phase composition} is then the statement that the generating function
of the whole process is the second phase's, composed over the state the first
phase has reached at $t_c$:
\begin{equation}
G^{\text{whole}}(x,t)=
\begin{cases}
G^{(1)}(x,t), & t<t_c,\\[4pt]
\displaystyle\sum_{k\ge0}\pr{X^{(1)}_{t_c}=k}
\bigl(G^{(2)}(x,t-t_c)\bigr)^{k}, & t\ge t_c,
\end{cases}
\label{var:eq:phasepgf}
\end{equation}
the $k$th power being available because the $k$ agents handed over found
independent lineages, and the weights being read off the first phase in the
usual way, $\pr{X^{(1)}_{t_c}=k}=(k!)^{-1}\partial_x^{\,k}G^{(1)}(0,t_c)$.

\begin{figure}[!htbp]
\centering
\begin{tikzpicture}[
  >={Stealth[length=1.9mm]},
  ph/.style={draw=ink,thick,rounded corners=2pt,minimum height=9mm,
             inner xsep=6pt,font=\small,align=center,text=ink},
  lb/.style={font=\footnotesize,inner sep=2pt},
  every path/.style={thick}
]
\node[ph,fill=type1!12] (p1) at (0,0)
  {phase 1: BDC\\[-1pt]\footnotesize $\lambda_1,\mu_1,\delta_1$};
\node[ph,fill=type2!12] (p2) at (6.2,0)
  {phase 2: BD\\[-1pt]\footnotesize $\lambda_2,\mu_2$};
\node[lb] (k) at (3.1,0.55) {$\Kb$ released};
\draw[->] (p1) -- node[below,lb]{$t_c$} (p2);
\node[lb,anchor=west] (est) at (8.2,0.45) {establishes, w.p.\ $1-p_0$};
\node[lb,anchor=west] (die) at (8.2,-0.45) {dies out, w.p.\ $p_0$};
\draw[->] (p2.east) -- (est.west);
\draw[->] (p2.east) -- (die.west);
\node[ph] (h0) at (0,2.4) {one infected host};
\node[ph] (h1) at (6.2,2.4) {two infected hosts};
\draw[->] (h0) -- node[above,lb]{w.p.\ $p=1-p_0$} (h1);
\draw[->,densely dashed] (h0.south) -- node[left,lb]{within-host} (p1.north);
\draw[->,densely dashed] (est.north) to[out=90,in=-30] (h1.south east);
\node[lb] at (3.1,3.35)
  {\itshape the host-to-host layer sits above the within-host process};
\end{tikzpicture}
\caption{The two-phase model and the branching layer above it. Within one
host the infection runs as a birth--death--catastrophe process until rupture
at $t_c$ releases a burst of size $\Kb$, and those $\Kb$ agents found
independent lineages of a slightly supercritical birth--death process, which
either establishes or dies out. Above it sits a Galton--Watson process with a
two-point offspring law: establishment sends the infection to two new hosts,
failure ends the line. The whole of the analysis below is the computation of
the one number $p_0$ joining the two levels.}
\label{var:fig:phase-schematic}
\end{figure}

\subsection{The host-to-host layer}
\label{var:sec:hosttohost}

Above the within-host process sits an epidemic. An infected individual either
goes on to infect exactly two others or recovers and infects none, and which
of the two happens is decided by whether the within-host infection
establishes. Writing $p$ for the probability that it does, the offspring
count $Z$ has
\begin{equation}
\pr{Z=2}=p,\qquad \pr{Z=0}=1-p,\qquad \ex{Z}=2p .
\label{var:eq:offspring}
\end{equation}

The two-point law is a modelling assumption and not a derived one. It is the
crudest available caricature of transmission, and it is adopted because it
makes $\ex{Z}$ a function of $p$ alone, $p$ being the one quantity the
within-host model supplies. Little rests on it: a general offspring law with
mean $m$ gives $\ex{Z}=mp$, so only the multiplier moves and nothing below
about how $\ex{Z}$ responds to the phase-1 parameters is an artefact of the
choice. What it does cost is the variance of $Z$, and with it every statement
about the epidemic beyond its mean --- the extinction probability of the
Galton--Watson process itself depends on the whole offspring law and not
merely on where it is centred.

\subsection{The extinction probability, derived}
\label{var:sec:extinction}

The within-host process fails if the phase-2 population dies out, and by the
law of total probability, conditioning on the burst size,
\begin{equation}
p_0=\sum_{k\ge0}\pr{\Kb=k}\lrb{\frac{\mu_2}{\lambda_2}}^{\!k}.
\label{var:eq:p0sum}
\end{equation}
The sum is worth naming: \cref{var:eq:p0sum} is the \emph{phase-1 burst-size
generating function evaluated at the phase-2 extinction probability},
$p_0=G_{\Kb}(\mu_2/\lambda_2)$. Both of its ingredients are proved elsewhere
in this thesis, and neither is re-derived here.

The first is the burst-size law. For a birth--death--catastrophe process the
load released at rupture is geometric on the larger root,
$\pr{\Kb=k}=(\delta_1/\lambda_1)a^{-k}$ for $k\ge1$, which is
\cref{dist:eq:burstlaw}, with $a$ and $b$ the roots of the characteristic
quadratic of \cref{bdc:lem:roots}. The second is the extinction probability
of a supercritical linear birth--death process from one founder,
$\mu_2/\lambda_2$, and hence $(\mu_2/\lambda_2)^k$ from $k$ independent
founders, which is \cref{m:sec:bd}.

The $k=0$ term is the mass on paths that never rupture at all: the cell's
occupants die out before the catastrophe can come, no load is released, and
no onward host is infected. That mass is $b$, and it is not a free
parameter.
\Cref{dist:eq:burstlaw} is stated unconditionally, including the paths that
never burst, and its normalisation is
$\sum_{k\ge1}(\delta_1/\lambda_1)a^{-k}=(\delta_1/\lambda_1)/(a-1)=1-b$, so
\begin{equation}
b=1-\frac{\delta_1/\lambda_1}{a-1},
\label{var:eq:bdetermined}
\end{equation}
which is also the smaller root of the same quadratic. The three constants
$a$, $b$ and $\delta_1/\lambda_1$ are determined by $(\lambda_1,\mu_1,\delta_1)$
together and only two of them may be chosen independently.

The sum in \cref{var:eq:p0sum} is geometric and closes. With
$x=\mu_2/(a\lambda_2)$,
\[
\sum_{k\ge1}\frac{\delta_1}{\lambda_1}\,a^{-k}
\lrb{\frac{\mu_2}{\lambda_2}}^{\!k}
=\frac{\delta_1}{\lambda_1}\cdot\frac{x}{1-x},
\]
so that
\begin{equation}
p_0=b+\frac{\delta_1}{\lambda_1}\cdot\frac{\mu_2}{a\lambda_2-\mu_2},
\qquad\text{convergent if and only if } a\lambda_2>\mu_2 ,
\label{var:eq:p0closed}
\end{equation}
and therefore
\begin{equation}
\ex{Z}=2\lrb{1-b-\frac{\delta_1}{\lambda_1}
\cdot\frac{\mu_2}{a\lambda_2-\mu_2}} .
\label{var:eq:EZ}
\end{equation}
The convergence condition costs nothing: the second phase is supercritical by
assumption, so $\lambda_2>\mu_2$, and $a>1$ by \cref{bdc:lem:roots}, so
$a\lambda_2>\lambda_2>\mu_2$ always.

What \cref{var:eq:p0closed} buys is a dependence of the epidemic on the
intracellular phase that is explicit rather than numerical. The first phase
appears in \cref{var:eq:EZ} only through $a$ and $\delta_1/\lambda_1$, and
both fall as $\delta_1$ falls, a cell less inclined to rupture holding its
occupants longer and releasing more of them. The second enters only through
the ratio $\mu_2/\lambda_2$, and the product $a\lambda_2$ in the denominator
is the whole of the interaction between them, and it is where the leverage
is: a larger burst multiplies the founding cohort of a near-critical process,
and since extinction from $k$ founders is the $k$th power of a number close
to one, multiplying $k$ is the most effective thing that can be done to it.

\subsection{Other compositions}
\label{var:sec:variants}

The composition above is the birth--death--catastrophe to birth--death case.
Three nearby variants share the same two-phase skeleton and differ only in
which parameter does the work: a first phase that is an ordinary birth--death
process running for a fixed time; the BDC case just analysed, in which
$\delta_1$ sets duration and handover together; and a first phase that simply
delivers a fixed count, the comparator against which the others are read.
A fourth variant does not share the machinery. It is a fixed handover into
the decaying death rate of \cref{var:rem:mudecay}, and with $\mu_2$ a
function of time the extinction probability of the second phase is no longer
$\mu_2/\lambda_2$ but the solution of the backward Riccati equation
\cref{m:eq:riccatiext}. It is named here and left there.

\begin{figure}[!htbp]
\centering
\includegraphics[width=0.72\textwidth]{fig12_EZ/fig12}
\caption{The host-to-host mean $\ex{Z}$ against the phase-1 catastrophe rate
$\delta_1$, from the closed form \cref{var:eq:EZ}, at
$(\lambda_1,\mu_1)=(1,0.2)$ and $(\lambda_2,\mu_2)=(1,0.9)$. The second phase
is barely supercritical, so a lineage founded by a single agent almost
certainly fails; what rescues the infection is arriving in numbers. A small
$\delta_1$ is a longer first phase and a larger release, and that is what
carries $\ex{Z}$ above one. The unit crossing is at $\delta_1\simeq0.061$.}
\label{var:fig:rough1}
\end{figure}

\begin{figure}[!htbp]
\centering
\includegraphics[width=0.96\textwidth]{fig13_EZ_decomp/fig13}
\caption{The same dependence, decomposed. Panel~(a): mean burst size against
$\delta_1$, a lower propensity to burst yielding a larger burst.
Panel~(b): $\ex{Z}$ against mean burst size, a larger founding cohort raising
the chance of surviving the near-critical extracellular phase.
Read together, $\delta_1$ acts on the epidemic entirely through the burst
size, which is what \cref{var:eq:EZ} makes explicit.}
\label{var:fig:rough23}
\end{figure}

\openqs{
\item \textbf{The decaying-rate variant.} The decaying-$\mu_2(t)$ model of
\cref{var:rem:mudecay} is named and not analysed. Taking it up needs the
inhomogeneous extinction probability of \cref{m:eq:riccatiext} in place of
$\mu_2/\lambda_2$.
\item \textbf{Against the threshold budget.} \Cref{path:sec:ypestis} reaches
the same organism and the same transition with different machinery, and
offers a threshold budget $\vartheta^{\ast}$ at which the intracellular
strategy stops paying. How $\ex{Z}$ here compares against that threshold ---
whether the $\delta_1$ that puts $\ex{Z}$ above one is the one that puts the
switch where $\vartheta^{\ast}$ puts it --- is the question that would join
the two accounts.
}

\FloatBarrier
